\lysh{14615}{4}

\begin{alist}
\item Παρατηρούμε ότι η εξίσωση είναι στην μορφή $\alpha x^2 + \beta x + \gamma = 0$ με
$\alpha = 1$, $\beta = -2\lambda$, $\gamma = \lambda^2 - 1$ και
$\Delta = (-2\lambda)^2 - 4 \cdot 1 \cdot (\lambda^2 - 1) = 4\lambda^2 - 4\lambda^2 + 4 = 4 > 0$, άρα η διακρίνουσα $\Delta$ είναι πάντα
θετική, ανεξάρτητα από την οποιαδήποτε τιμή της παραμέτρου $\lambda$.

\item 
1ος τρόπος:
$\rho_{1,2} = \dfrac{-\beta \pm \sqrt{\Delta}}{2\alpha} = \dfrac{-(-2\lambda) \pm \sqrt{4}}{2 \cdot 1} = \dfrac{2\lambda \pm 2}{2} = \dfrac{2(\lambda \pm 1)}{2} = \lambda \pm 1$. Ώστε $\rho_1 = \lambda - 1$, $\rho_2 = \lambda + 1$.

2ος τρόπος:
Η εξίσωση γράφεται $(x - \lambda)^2 - 1^2 = 0$ και ισοδύναμα έχουμε
$(x - \lambda - 1)(x - \lambda + 1) = 0$ άρα
$x - \lambda - 1 = 0$ ή $x - \lambda + 1 = 0$. Ώστε $x = \lambda + 1$ ή $x = \lambda - 1$.
Έτσι $\rho_1 = \lambda - 1$, $\rho_2 = \lambda + 1$ αφού $\rho_1 < \rho_2$.

\item Πρέπει $|\rho_2 - (-\rho_1)| \ge 8 \Leftrightarrow |\rho_1 + \rho_2| \ge 8$. Αν έχουμε βρει τις ρίζες τότε $\rho_1 + \rho_2 = 2\lambda$.
Χωρίς να βρούμε τις ρίζες, από τύπους Vieta, $\rho_1 + \rho_2 = -\dfrac{\beta}{\alpha} = \dfrac{-(-2\lambda)}{1} = 2\lambda$.

Άρα, πρέπει $|2\lambda| \ge 8 \Leftrightarrow 2|\lambda| \ge 8 \Leftrightarrow |\lambda| \ge 4 \Leftrightarrow \lambda \le -4$ ή $\lambda \ge 4$.

\item Έστω το τριώνυμο $f(x) = 1 \cdot x^2 - 2\lambda x + \lambda^2 - 1$ τις ρίζες του οποίου έχουμε βρει στο $\beta)$ ερώτημα.
Παρατηρούμε ότι ο αριθμός $k^2 - 2\lambda k + \lambda^2 - 1 = f(k) < 0$, σχέση που προκύπτει από τον
παρακάτω πίνακα προσήμου του $f(x)$ για τον οποίο γνωρίζουμε ότι ανάμεσα στις ρίζες το τριώνυμο είναι ετερόσημο του $\alpha = 1 > 0$, συνεπώς για $k$ μεταξύ των ριζών $\rho_1, \rho_2$, είναι $f(k) < 0$.
Άρα ο $k$ βρίσκεται μεταξύ των ριζών $\rho_1$ και $\rho_2$.
\end{alist}

